% ==========================================================================
% GENERATED by scripts/split_paper.py. Do not edit this file.
%
% Source: imported_papers/grokking/iclr2027_conference_v2.tex
% Part 4: Results: Credit allocation changes circuit concentration
%
% Edits belong upstream, in the published paper. To change what the
% chapter includes, change imports/grokking.toml and re-run
% the splitter.
% ==========================================================================
\section{Results: Credit allocation changes circuit concentration}
\label{sec:select}

\paragraph{Measuring circuit concentration.} We ask how a learned computation is spread across the neurons of a layer in the networks that grok. On modular addition, algorithm checks for all the tested learning rules that grok show the same periodic Fourier computation and not the alternative computation reported by \citet{zhong2023clock} (Appendix~\ref{app:algo}). The checks cover the input frequencies used, a trigonometric fit of the logits, and agreement across seeds and rules; for the reference and the equalized network, the fitted trigonometric function alone classifies every held-out pair correctly.

We measure circuit concentration from the norms of the weights that connect each hidden neuron to the next layer, since, at a given activation, a neuron with small outgoing weights has only a small effect on the output, whatever it computes. For the $n$ neurons in the measured layer, let $a_j = \|\mW^{\mathrm{out}}_{:,j}\|_2$ be the norm of the outgoing weights of neuron $j$. Their \emph{participation ratio} is
\begin{equation}
\mathrm{PR} \;=\; \frac{\big(\sum_j a_j\big)^2}{n \sum_j a_j^2} \;\in\; [1/n, 1],
\label{eq:pr}
\end{equation}
for a nonzero vector of norms. If all norms are equal, $\mathrm{PR} = 1$, and if only one is nonzero, $\mathrm{PR} = 1/n$ (Proposition~\ref{prop:pr}, and all propositions are stated and proved in Appendices~\ref{app:theory} and~\ref{app:generalization}). Applied to the outgoing norms, as here, we call it the \emph{circuit participation}, and $n \cdot \mathrm{PR}$ the \emph{effective neuron count} (a layer with effective count $154$ has the participation ratio of a layer in which $154$ neurons carry equal outgoing norms and the rest none). We call a neuron \emph{nearly unused} if its outgoing norm is below a tenth of the layer mean. Both are computed from the weights, and Section~\ref{sec:rules} confirms them with two measures tied to the computation. Applied to the norms of the learning signal that the neurons receive (Section~\ref{sec:law}), we call it the \emph{credit participation}.

Backpropagation with Adam produces a concentrated circuit with approximately $154$ effective neurons out of $256$ in the first hidden layer, and about $84$ of the $256$ neurons are nearly unused. In contrast, DFC groks with a delay of $520$ to $610$ epochs (Appendix~\ref{app:repro}) and produces a distributed circuit with approximately $235$ effective neurons at nearly perfect test accuracy.

\subsection{Equalizing credit across neurons}
\label{sec:rules}

Backpropagation and DFC differ in several respects, so we modify backpropagation itself and keep everything else fixed.

Let $\vg_j$ be the gradient of neuron $j$'s incoming weights. Its norm $\|\vg_j\|_2$ is the neuron's credit magnitude. The mean credit magnitude over the $n$ neurons of the layer is $\bar{g} = \frac{1}{n}\sum_{k=1}^{n} \|\vg_k\|_2$. The intervention moves each magnitude toward the mean with strength $\beta \in [0,1]$,
\begin{equation}
\widetilde{\vg}_j \;=\; \frac{\vg_j}{\|\vg_j\|_2}\Big( (1-\beta)\, \|\vg_j\|_2 + \beta\, \bar{g} \Big)
\quad \text{for } \|\vg_j\|_2 > 0,
\label{eq:knob}
\end{equation}
and leaves zero gradients at zero. At $\beta = 0$ the gradient is unchanged and training is ordinary backpropagation with Adam. At $\beta = 1$ all nonzero incoming-weight gradients have the same norm. We call this \emph{full credit equalization}. The direction of every neuron's incoming-weight gradient is preserved, so before it enters Adam the rescaled gradient is still a descent direction for the loss (Proposition~\ref{prop:descent}, Appendix~\ref{app:theory}). The rescaling acts only on the incoming-weight gradients of the first hidden layer. Circuit participation is measured from the outgoing weights of the same neurons, which are the incoming weights of the second hidden layer and are not rescaled.

\begin{figure}[t]
\begin{center}\includegraphics[width=\linewidth]{figures/fig2_switch.pdf}\end{center}
\caption{Credit allocation sets the circuit and reorganizes it after grokking (modular addition, $p=31$, three seeds, bands one standard deviation). \textbf{(a)} Circuit participation and effective neuron count of networks trained from scratch at equalization strength $\beta$. Training cross-entropy ($0.040$), margin and parameter norm stay close to the reference values. \textbf{(b)} A network grokked with the reference rule and one grokked with credit equalization, each copied at epoch $20{,}000$ and continued for $30{,}000$ epochs with its own rule or with the other rule. Test accuracy remains above $0.95$ throughout.}
\label{fig:steer}
\end{figure}

Increasing $\beta$ from $0$ to $1$ raises circuit participation from $0.60$ to $0.99$, that is, from approximately $154$ to $254$ effective neurons (Figure~\ref{fig:steer}a). From $\beta = 0.25$ on no neuron is nearly unused, yet participation keeps rising, from $0.93$ to $0.99$ (Appendix~\ref{app:adapt}). Training loss, margin (the mean difference between the correct logit and the largest incorrect logit) and parameter norm stay within a few percent of the reference values, and mean test accuracy decreases from $0.998$ to $0.976$. The gap widens with width: at width $2048$ the reference uses $332$ effective neurons and the equalized network $2046$ (Appendix~\ref{app:arch}). The circuit concentrates during the grokking delay. In the reference, participation is $0.997$ when the training set is first fitted, has fallen to between $0.59$ and $0.67$ when test accuracy reaches $0.95$, and changes by at most $0.05$ afterwards (five seeds, Appendix~\ref{app:timing}).

Two measures that use each neuron's activity or its effect on the output also increase. Over eight seeds, participation computed from outgoing norms weighted by mean absolute activation rises from $0.34$ to $0.81$. Participation computed from the absolute mean change in the correct-class logit after removing each neuron rises from $0.42$ to $0.80$ (Appendix~\ref{app:functional}). The activation-weighted measure is invariant to compensating rescalings of a ReLU neuron's weights.

We next ask whether the intervention can redistribute a circuit that has already been learned. We let the reference network grok, then switch to full credit equalization and continue training on the same task. Outgoing-weight participation rises from $0.60$ to above $0.95$ within $300$ epochs and then to $0.99$, while test accuracy remains above $0.98$ (Figure~\ref{fig:steer}b). The effect of the neurons on the output spreads as well, from ablation participation $0.40$ to $0.72$ over five seeds (Appendix~\ref{app:adapt}). A control that continues with the reference rule stays at $0.60$, and in separate runs both rules hold their participation over $200{,}000$ epochs (Appendix~\ref{app:long}).

Controls in Appendix~\ref{app:credit_controls} exclude alternative explanations. Constraining the incoming weight norms instead of the credit gives participation $0.66$, and a grid over learning rate and weight decay stays between $0.53$ and $0.63$. Dropout also distributes the circuit, but it changes the objective and raises margin and norm (Appendix~\ref{app:baselines}). In a simplified linear model under gradient descent, each unit grows in proportion to its size without equalization and by the same amount with it, so initial differences persist in the first case and fade in the second (Proposition~\ref{prop:select}, Appendix~\ref{app:theory}).

\paragraph{A per-neuron condition.} A neuron stays in use when its credit is large relative to the non-gradient term of its own update. Adam is nearly invariant to a constant rescaling of its input, so rescaling a neuron's gradient changes its size relative to the other terms of its update, in our reference the coupled decay term added to the gradient before Adam. Rescaling only this term per neuron, with the gradient unchanged, reproduces the redistribution, and equalizing the gradient while holding each neuron's ratio of gradient to decay leaves the circuit concentrated, on the activation-weighted and ablation measures as well (Table~\ref{tab:pathways}, Appendix~\ref{app:decay}). The same condition holds for two other terms that do not scale with the gradient: Adam's $\epsilon$, which dominates the denominator of its update when AdamW decays a transformer's normalization parameters (Appendix~\ref{app:transformer}), and the decay of a trainable per-neuron gain and of the outgoing weights (Appendix~\ref{app:normlayers}). Lion and SGD with coupled decay follow it as well (Appendix~\ref{app:baselines}). Where no non-gradient term dominates, as under AdamW without decay on normalization parameters, the reference is already broadly distributed, with participation $0.91$.

\begin{table}[t]
\caption{Rescaling only the non-gradient term of each neuron's update reproduces credit equalization in all three cases. Equalizing the gradients while holding each neuron's ratio of credit to that term fixed keeps the circuit concentrated for the decay and the gain and removes about half of the effect for $\epsilon$ (circuit participation, five seeds per cell). For the gain, equalization and rescaling act on the gain and the outgoing weights as well as the incoming rows; equalizing the incoming rows alone gives $0.364$. Update equations and details in Appendices~\ref{app:decay}, \ref{app:transformer} and~\ref{app:normlayers}.}
\label{tab:pathways}
\begin{center}
\footnotesize
\setlength{\tabcolsep}{4pt}
\begin{tabular}{llcccc}
\toprule
Non-gradient term & Network & Reference & Equalized & Term rescaled & Ratio held \\
\midrule
Coupled decay & MLP, Adam with L2 & $0.606$ & $0.993$ & $0.993$ & $0.574$ \\
Adam's $\epsilon$ & Transformer, AdamW & $0.609$ & $0.994$ & $0.981$ & $0.818$ \\
Decay of the gain & MLP, layer normalization & $0.366$ & $0.996$ & $0.997$ & $0.343$ \\
\bottomrule
\end{tabular}
\end{center}
\end{table}

\subsection{Which parameters are normalized together?}
\label{sec:law}

Full credit equalization is itself a normalization of the gradient with one factor per neuron. Optimizers normalize the learning signal at other scales: Adam normalizes each weight separately, LARS rescales a whole layer by one factor, and Nero normalizes the gradient of the weights into each neuron. Before comparing optimizers, we ask in a controlled setting whether normalizing several neurons together has the same effect as normalizing each one alone. We divide the $n$ neurons of the first hidden layer into groups of $m$ neurons and rescale all gradients in a group by one common factor, chosen so that every group has the same mean credit,
\begin{equation}
\widetilde{\vg}_j \;=\; \vg_j\, \frac{\bar{g}}{\bar{g}_{B(j)}},
\qquad
\bar{g}_{B(j)} \;=\; \frac{1}{m}\sum_{k \in B(j)} \|\vg_k\|_2,
\label{eq:grain}
\end{equation}
where $B(j)$ is the group containing neuron $j$ and $\bar{g}$ is the layer mean from Equation~\ref{eq:knob}. With $m = 1$ every neuron is rescaled separately and the rule is full credit equalization ($\beta = 1$). With $m = n$ the factor equals one and the gradient is unchanged.

For groups of eight or more neurons, circuit participation stays within about $0.02$ of the value $0.60$ of the reference, backpropagation with Adam. It rises to $0.70$ for groups of four, $0.75$ for pairs and $0.99$ for single neurons (Figure~\ref{fig:law}b). Most of the redistribution occurs only when each neuron receives its own factor. This equalizes the mean credit between groups but keeps the ratios inside a group: a neuron that received twice the credit of its neighbor still does after the rescaling. In a simplified model, a linear network in which every unit has one incoming and one outgoing scalar weight, these within-group differences persist throughout training and keep the circuit concentrated. Only single-neuron groups leave none to preserve (Proposition~\ref{prop:grain} in Appendix~\ref{app:theory}, sweep details in Appendix~\ref{app:grain}).


\paragraph{Comparison with optimizers.} If per-neuron normalization is what distributes the circuit, optimizers should sort accordingly. We train with the eight optimizers of Table~\ref{tab:ruleset} that replace Adam. LARS, RMSProp, LAMB, Adafactor, Lion, Muon and Shampoo normalize groups other than a single neuron or change update directions, and all seven produce concentrated circuits, with participation between $0.470$ and $0.613$. All seven run with the coupled weight decay of the reference. Nero, which uses no weight decay, produces at its default learning rate of $10^{-2}$ a distributed circuit with participation $0.986$ at test accuracy $0.985$ (Table~\ref{tab:zoo}). Nero, which normalizes credit per neuron, and DFC, whose credit is nearly equal (credit participation $0.94$), are cases of equal credit that arise without our intervention, and both produce distributed circuits.

For the reference, DFC and the equalization sweeps, credit participation lies close to circuit participation (Figure~\ref{fig:law}a). The condition also explains why equal update norms do not suffice. Lion gives every neuron an update of the same norm, and Adam-mini \citep{zhang2025adammini} normalizes each neuron's update with one second-moment estimate, yet both keep the circuit concentrated (Adam-mini $0.586$ against $0.594$ for Adam, Appendix~\ref{app:grain}). Both normalize the coupled decay together with the gradient, so each neuron's ratio of credit to decay is unchanged.

\begin{figure}[tb]
\begin{center}\includegraphics[width=\linewidth]{figures/fig_law.pdf}\end{center}
\caption{Circuit participation under different transformations of the learning signal. \textbf{(a)} Circuit participation against credit participation for the reference, DFC, credit equalization at $\beta = 0.25$ and $\beta = 1$, the group-size sweep on four tasks and seven optimizers. Credit participation is measured on the signal listed in Table~\ref{tab:creditconv}. \textbf{(b)} Circuit participation against the number of neurons that share one normalization factor, on modular addition at three moduli and on $S_5$. Error bars show one standard deviation across three seeds.}
\label{fig:law}
\end{figure}

\subsection{Which learning rules grok}
\label{sec:grokrules}
On modular addition, backpropagation with the optimizers in Table~\ref{tab:ruleset}, credit equalization, DFC and the two constructed rules with a fresh or a fading orthogonal component grok, whereas feedback alignment, direct feedback alignment, sign-SGD and the rule with a fixed orthogonal component fit the training set and stay near zero test accuracy over $15{,}000$ epochs. This split matches the weight-decay account of grokking \citep{nanda2023progress,varma2023circuit}. Once the training data are fitted, weight decay removes weight along the directions in which small parameter changes leave the training outputs unchanged, and a learning signal that is not a loss gradient can hold weight there. The norm $q$ of the weight in these directions (Appendix~\ref{app:generalization}) separates the two outcomes at the end of training without overlap across nine configurations with three seeds each (Figure~\ref{fig:main}a). It shrinks to means between $5.3$ and $8.5$ in the rules that grok and stays between $26$ and $135$ in the rules that memorize. Within a rule, $q$ does not predict the outcome (Appendix~\ref{app:generalization}).
